\section{主观题}
\begin{enumerate}
  \item 一个案例，带 4 - 6 个小问题，问题相对独立。
  \item 考能力为主。阅读案情，\bs{识别考点}，分析问题，搜索法条，下结论。
  \item \zd{不要求必须明确说出依据的法律规定的来源和具体的条文标号}。
  \item \zd{核心有用的信息}。
  \item 考点分布。
        \begin{enumerate}
          \item 考流程。围绕审判为中心。一审。
          \item 证据。
          \item 新法。
        \end{enumerate}
  \item 案情设计。
        \begin{enumerate}
          \item 不关心罪名，只为考点服务。
          \item 人物线、时间线（问题设问按时间顺序）。权力、权利。
        \end{enumerate}
  \item 题型。
        \begin{enumerate}
          \item 救济类。考人。救济方式 + 法律依据/理由。
          \item 方案类。考机关。有几个分几段，方案写具体，结合案情，不需要论证。\al{指令 S 中院重审}。
          \item 是非评价类。短句评价，长句说理。\al{是否妥当}。
          \item 纠错。好多个错。错在哪 + 理由。不流行。
          \item 简答论述。\al{拘传、传唤区别}。
          \item 开放式。三段论。\al{是否分案}。
        \end{enumerate}
  \item 要求。
        \begin{enumerate}
          \item 最多读 8 分钟。
          \item 问题反向定位案情。
        \end{enumerate}
  \item 采分点。结论是否正确、论证说理、法言法语。
        \begin{enumerate}
          \item 答题给出结论和理由要明确直接。\zd{按点采分}。
          \item 用最明确的语言给出\bs{简短直接}的答案。
          \item \zd{分点}阐述。
          \item 专业术语 \hz 法律规范的规定来直接明确作答，\bs{不需要过多展开论述}。
        \end{enumerate}
  \item \zd{知识点}。
        \begin{enumerate}
          \item 每个诉讼阶段的任务、各诉讼主体的权利义务。
          \item 阶段推进的条件、标准。\al{侦查终结、移送审查起诉、提起公诉、二审启动、再审启动}。
        \end{enumerate}
  \item \bs{法条定位}。
        \begin{enumerate}
          \item 宏观找诉讼。
          \item 微观找各妈。高法解释（\zd{刑事诉讼法}）、高检规则（\zd{诉讼规则}）、公安部规定（\zd{公安机关}）。
          \item 专项找专规。
                \begin{enumerate}
                  \item \zd{认罪认罚}。两高三部、\zd{量刑建议}。
                  \item 排非有两个，用 17 年那个。\zd{严格排除}。
                  \item 规范量刑。\zd{量刑程序}。
                  \item \zd{值班律师}。
                  \item \zd{电子数据}。
                  \item 职务犯罪。\zd{监察法}。
                \end{enumerate}
          \item \bs{先阶段 \yjt 再具体名词}。\al{审查起诉阶段，发现强疗}。
          \item 新增。\bs{信息网络、取保候审}。
          \item 宏观 \yjt 微观。考点 \yjt 文件 \yjt 章节 \yjt 条文
        \end{enumerate}
\end{enumerate}
